% Policy layer (review sections 4 and 9.3).
\textbf{Trusted policy layer.} The VLM proposes tool calls; it never decides whether they run. Every call, whether
it comes from the model, the rule-based fallback or a \texttt{/}command, passes a deterministic policy layer
(the ``action-selector'' and ``dual-LLM'' patterns~\cite{beurerkellner2025patterns}; complete mediation as
recommended for excessive agency~\cite{owasp2025llm}): (i)~\emph{roles}: viewers may read, operators may change
alert rules and the counting line, and only administrators may remove \emph{all} alerts; the role comes from the
authenticated token, never from the chat text; (ii)~\emph{validation}: types, finite numbers, per-metric ranges
(a threshold of a million people is a disabled alert in disguise), line length and position; (iii)~an
\emph{intent gate}: a state change proposed by the model must be requested in the user's own words of this turn
(a change verb and its object), so text that reached the model from the frame, a scene report or a tool result
cannot trigger one; (iv)~\emph{rate limits} per principal on changes and chat requests, and refusal when the VLM
queue is full; (v)~\emph{confirmation}: removing or weakening an alert from chat waits for \texttt{/confirm
<id>} from the same principal; and (vi)~a \emph{hash-chained audit log} of every decision
($h_n = \mathrm{SHA256}(h_{n-1}\,\|\,\text{record}_n)$~\cite{schneier1999audit}), which detects edited, deleted or
reordered entries. In addition, stage~1 is \emph{isolated} from VLM-written text: it never sees scene reports,
earlier VLM answers (only the system-written footer of each turn) or the free-text fields of events, so words read
from the camera view cannot reach the only component that can change settings. Events, series, reports and chat
are deleted after a retention period (72\,h by default); the audit log is kept.
